% =============================================================================
%  «ТЕОРИЯ СИСТЕМ» — ствол серии «Архитектура Рассуждения»
%  Блок 1. ВХОД: от метафизики к теории.  Разделы 1.1–1.2 (черновик прозы).
%  Стиль: «выводим и показываем» — логические формулы в тексте; в каждом разделе
%  сначала ВЫВОД/ПОКАЗ, затем РЕЗУЛЬТАТЫ; репо-якоря — в приложении, не в сносках.
%  Позитивная онтология: начинаем с единицы, не с нуля.
%  Самокомпилируемо: pdflatex blok-1-vhod.tex (см. память latex-build-windows).
% =============================================================================
\documentclass[12pt,a4paper]{article}
\usepackage{cmap}            % ToUnicode -> кириллица копируема/ищется
\usepackage[T2A]{fontenc}
\usepackage[utf8]{inputenc}
\usepackage[russian]{babel}
\usepackage{paratype}        % PT Serif/Mono — масштабируемый, копируемая кириллица
\usepackage{amsmath,amssymb,amsthm}
\usepackage{listings}
\usepackage[expansion=false]{microtype}
\usepackage[margin=28mm]{geometry}

% --- листинг для Coq-кода (ASCII; без подсветки ключевых слов, ради надёжности) ---
\lstset{
  basicstyle=\ttfamily\small,
  frame=single,
  framesep=6pt,
  columns=fullflexible,
  keepspaces=true,
  breaklines=true,
  xleftmargin=10pt,
  aboveskip=10pt, belowskip=10pt
}

% --- лёгкий маркер фазы «результаты» (стиль книги: показ -> затем результаты) ---
\newcommand{\rezultat}{\medskip\noindent\emph{Что отсюда следует.}\ }

\title{\textbf{Теория Систем}\\[2mm]\large Блок 1. Вход: от метафизики к теории}
\author{}
\date{}

\begin{document}
\maketitle
\thispagestyle{empty}

\noindent\small\textit{Черновик: разделы 1.1--1.5 (Блок~1). Глубинная метафизика (Логика как
структура бытия, Закон Достаточного Основания как её действие, Источник, Свидетель,
каскад систем) принадлежит Тому~I «Метафизика» и здесь принимается как опора, не
переоткрывается. Машинные якоря — в Приложении в конце блока, не в сносках.}\normalsize

\bigskip

% =============================================================================
\section*{1.1\quad Первичный факт: $A=\exists$}
% =============================================================================

Книга начинается не с определения системы и не с аксиомы, а с факта, отрицание
которого само себя опровергает.

Существовать --- значит существовать \emph{определённо}, то есть быть отличимым от
того, чем не являешься. Запишем первичный факт символически:
\[
  A \;=\; \exists \qquad\text{(нечто есть).}
\]
Оговорим саму запись. Знак $\exists$ здесь \emph{не} квантор внутри формального
равенства, и $A=\exists$ не есть равенство некоторого терма с квантором. Это краткое
символическое указание на первичный факт: \emph{нечто есть, существование дано}.

Отрицание этого факта невозможно без впадения в противоречие. Пусть кто-то
утверждает $\neg\exists$ --- «ничего нет». Само это утверждение требует утверждающего,
акта утверждения и различения «нет» от «да»; тем самым акт отрицания
\emph{инстанцирует} то, что отрицает:
\[
  \neg\exists \;\rightarrow\; \exists,
\]
а значит $\exists$. У отрицания нет основания, которое не предполагало бы
существование. Это работа Закона Достаточного Основания на самом первичном уровне:
положение утверждает себя, ибо для его отрицания нет основания, не порождающего
противоречия.

Заметим \emph{границу записи}. Сам факт $A=\exists$ не может быть первой строкой
формальной системы: всякая запись предполагает уже наличный язык, типы, средства
различения. Поэтому формализация начинается не с факта существования, а с его первого
\emph{следствия} --- с акта различения, к которому мы перейдём в~§1.2.

\rezultat
Существование \emph{обнаруживается}, а не постулируется: это не первая аксиома
(формальной системы в этой точке ещё нет) и не философский выбор, а факт,
установленный анализом работы Логики на самой себе. Отсюда --- позиция, которую книга
держит всюду: \textbf{мы начинаем с наличия, не с пустоты.} Это позитивный старт.

Подчеркнём, однако, чего здесь ещё \emph{нет}: здесь нет числа. Сказано лишь, что
нечто есть, --- но не «одно». \emph{Единица} как первое количество рождается не из
факта существования, а из акта различения; к ней мы переходим сейчас.

% =============================================================================
\section*{1.2\quad Акт различения: рождение единицы}
% =============================================================================

Первое, что записывается, --- акт различения. Существовать определённо значит быть
отделённым: провести границу между $A$ и $\neg A$. У этого акта жёсткая структура --- в нём
\emph{две стороны} и \emph{два условия}, скрепляющие их:
\begin{itemize}
\item две стороны --- \emph{выделяемое} $A$ и \emph{фон} $\neg A$;
\item \emph{условие непротиворечия} (Закон~L2): стороны не пересекаются, $\;\neg\,(A \wedge \neg A)$;
\item \emph{условие исчерпания} (Закон~L3): третьего не дано, $\;A \vee \neg A$.
\end{itemize}
Больше в различении ничего нет: две стороны, удержанные вместе непересечением и исчерпанием.
(В машинной проверке этой структуре отвечает одна запись из четырёх полей; её точная форма ---
в Приложении, и для понимания она здесь не нужна.)

В одном этом акте \emph{неразрывно} присутствуют три грани. \emph{Что} различается ---
субстрат $A$ --- это Elements. \emph{Как} различается --- сам разрез $A\mid\neg A$ ---
это Roles. \emph{Почему} так --- законы, управляющие разрезом --- это Rules. Это не три
части, которые можно собрать порознь, а три грани \emph{одного} акта: невозможно
различение без всех трёх сразу. Здесь --- зерно триады E/R/R, к которой цепь придёт
в~§1.5.

Обе стороны рождаются одним актом \emph{одновременно}: для различения, построенного по
предложению $P$, выделяемая сторона есть $P$, а фон есть $\neg P$, и обе фиксируются
единым построением --- ни одна не предшествует другой.

\medskip
Теперь --- количественное начало. Различения можно \emph{считать}: сопоставим набору
различений его число. Пустому набору --- когда не совершено ни одного акта --- отвечает
число нуль:
\[
  0 \;\longleftrightarrow\; [\,]\ \ (\text{пустой набор}) \;\longleftrightarrow\;
  \textbf{отсутствие.}
\]
Одному акту различения отвечает число один:
\[
  1 \;\longleftrightarrow\; [\,d\,]\ \ (\text{одно различение}) \;\longleftrightarrow\;
  \textbf{первая определённость, единица.}
\]
Дальнейший счёт есть добавление одного различения за другим: от $n$ к $n+1$ --- ещё
один акт. Так натуральный ряд не постулируется, а \emph{порождается} актом различения,
шаг за шагом.

Здесь же видна асимметрия, на которой держится начало с единицы. Условие непротиворечия
$\neg(A\wedge\neg A)$ читается направленно: \emph{удерживая выделяемую сторону, мы тем самым
отрицаем фон}. Положительное \emph{само\-заземлено} --- оно обосновывает отрицание своего
фона; отрицательное (а с ним отсутствие и нуль) --- \emph{производно}. Это Закон Достаточного
Основания в форме самозаземления.

\rezultat
Единица не вводится постулатом «пусть дано~$1$». Она \emph{есть} число одного акта
различения --- первой актуальной определённости. Поэтому счёт начинается с единицы:
первая актуальность --- одно различение. Нуль же не первичен: он есть число
\emph{пустого} набора, имя отсутствия. Отсюда натуральный ряд начинается с единицы,
\[
  \mathbb{N} \;=\; \{\,1,\,2,\,3,\,\dots\,\},
\]
а то, что в средстве машинной проверки тип натуральных чисел отсчитывается от нуля, ---
техническое соглашение инструмента, не онтологическая позиция работы.

Так формулируется позитивная онтология уже \emph{точно}: мы начинаем с единицы --- с
одного акта, с положительного, --- а отсутствие и нуль производны. Это первая из двух
сквозных нот книги. Вторая --- \emph{граница} --- зазвучит, когда от законов мы перейдём
к принципам и встретим конечную актуальность (§1.4): бесконечность есть свойство
процесса, а не завершённого объекта. Книга открывается единицей и закроется границей.

% =============================================================================
\section*{1.3\quad Пять законов как теоремы различения}
% =============================================================================

Акт различения уже принёс два закона: непротиворечие $\neg(A\wedge\neg A)$ и исчерпание
$A\vee\neg A$ --- два условия, скрепившие стороны (§1.2). Соберём теперь полную пятёрку и
уточним её статус.

Законы Логики не вводятся здесь как новые допущения. Они суть свойства самой Логики,
со\-данные с бытием в том же первичном факте; в записи различения мы их не учреждаем, а
\emph{прослеживаем} как теоремы --- то, что уже работало в акте, выходит на поверхность
доказанным.

Пять законов читаются на одном разрезе $A\mid\neg A$:
\begin{itemize}
\item \emph{Тождество} (L1): выделенное есть оно само, $A$ есть $A$, --- сторона устойчива,
пока держится разрез.
\item \emph{Непротиворечие} (L2): стороны не пересекаются, $\neg(A\wedge\neg A)$.
\item \emph{Исключённое третье} (L3): третьего не дано, $A\vee\neg A$.
\item \emph{Достаточное основание} (L4): у разреза есть основание; в самой сжатой форме ---
самозаземление положительного, $\mathit{positive}\to\neg\,\mathit{negative}$ (§1.2).
\item \emph{Порядок} (L5): различения встают на уровни --- выделенное и его элементы не на
одной ступени, и эта иерархия не имеет конца внутрь себя.
\end{itemize}

Нужна оговорка о статусе. Четыре закона из пяти --- теоремы: они доказываются из устройства
акта. Один --- \emph{исключённое третье} (L3) --- принимается как несущая опора: это
\emph{единственная аксиома} всей ветви. Различение требует, чтобы сторона была либо выделена,
либо нет, без зазора; этот «без зазора» и есть L3, и обосновать его, не используя, нельзя.
Книга держит его открыто --- как одну принятую опору, а не прячет среди выводов.

О двойной нумерации. У каждого закона два имени: \emph{онтологическое} (Тождество,
Непротиворечие, …) --- чем закон является в бытии --- и \emph{формальное} (L1–L5) --- его метка
в порядке вывода. Это одно и то же, прочитанное с двух сторон: со стороны сути и со стороны
доказательства.

\rezultat
Пятёрка законов \emph{выведена} из одного акта различения, а не приставлена к нему.
\textbf{Логика не прилагается к различению извне --- она и есть его внутреннее устройство.}
Несущая опора ровно одна (L3); остальные четыре --- следствия. Тем самым «законы логики»
перестают быть списком данных правил и становятся \emph{прослеженной} структурой первого акта.

% =============================================================================
\section*{1.4\quad Четыре принципа существования}
% =============================================================================

Законы говорят, как устроен разрез. Принципы говорят, как устроено \emph{существование}
того, что разрезом выделено. Их четыре; три --- теоремы из пяти законов, четвёртый требует
отдельной честности.

\begin{itemize}
\item \emph{P1 --- нет само\-членства.} Выделенное не содержит само себя как собственный
элемент: его уровень и уровень его элементов различны (L5). Поэтому «множество всех множеств,
не содержащих себя» не \emph{строится} --- у него нет хорошей формы. Парадокс Рассела здесь не
разрешается задним числом; он \emph{невозможен}.
\item \emph{P2 --- критерий прежде системы.} Чтобы собрать нечто в систему, нужен признак
принадлежности, и он предшествует системе (L5): сначала «по чему собрано», затем «что
собрано». Этим снят круг «система через элементы, элементы через систему».
\item \emph{P3 --- тождество по содержанию.} Две вещи, неразличимые по всем признакам, суть
одна (L1+L4): различие требует основания-различителя; нет различителя --- нет и двух.
\item \emph{P4 --- конечная актуальность.} Актуально дано всегда конечное; бесконечность есть
свойство \emph{процесса}, а не завершённого объекта.
\end{itemize}

О статусе P4 --- честно. Первые три принципа доказываются из законов. Четвёртый \emph{не} есть
чистый вывод: он обеспечен \emph{выбором конструкции}. Процесс мы берём как последовательность
приближений, где каждая ступень --- рациональное число $n/d$, заданное конечно (по
секвенциальности L5); конструктивный переход «из возможности --- в актуальное» вынесен в
отдельную опору \emph{свидетеля} (названа в Приложении). При таком выборе актуально-бесконечное
не возникает: каждая ступень конечна, а бесконечность остаётся \emph{горизонтом} процесса ---
тем, к чему он приближается, не достигая.

\rezultat
Существование структурировано четырьмя принципами; три --- теоремы законов, а P1 заодно делает
классические парадоксы невозможными, а не лечит их. Но именно на P4 \emph{впервые звучит вторая
нота книги}: \textbf{актуальное конечно, бесконечность --- процесс, не вещь.} Это та граница,
которая в Блоке~7 окажется единственной несущей стеной теории; здесь она лишь подаёт первый
звон --- и подаёт честно, как выбор конструкции, а не как выведенную теорему.

% =============================================================================
\section*{1.5\quad E/R/R: конец цепи}
% =============================================================================

Цепь замыкается. В §1.2 мы видели в одном акте различения три неразрывные грани --- \emph{что}
(Elements), \emph{как} (Roles), \emph{почему} (Rules). Пройдя законы и принципы, можно назвать
получившееся его собственным именем: всякая система есть триада \textbf{E/R/R}.
\begin{itemize}
\item \textbf{Elements} --- то, из чего система; её наполнение.
\item \textbf{Roles} --- то, как элементы соотнесены; разрезы и связи.
\item \textbf{Rules} --- то, по каким законам соотнесение держится; конституция.
\end{itemize}

Это не три части, собираемые порознь, а три грани одного: деревяшка без роли и правила не есть
молоток; правило без носителя и роли не есть система. Порядок прочтения и порядок порождения
здесь \emph{обратны}. Познавая, мы идём от наличного наполнения к ролям и правилам:
Elements~$\to$~Roles~$\to$~Rules. Порождается же система в обратную сторону: правило задаёт
возможные роли, роли --- что вообще может стать элементом: Rules~$\to$~Roles~$\to$~Elements.
Эту асимметрию --- первичность правил --- Блок~2 докажет как теорему.

\rezultat
Триада E/R/R --- \emph{не постулат, а конец дедуктивной цепи}:
\[
  A=\exists \;\to\; \text{различение} \;\to\; \text{пять законов} \;\to\;
  \text{четыре принципа} \;\to\; \text{E/R/R}.
\]
Система перестаёт быть исходным понятием и становится \emph{выведенным} объектом. Цепь дала
объект --- и со следующего блока мы разбираем его изнутри: как он устроен, как складывается с
другими, как живёт во времени и где его единственная граница.

% =============================================================================
\appendix
\section*{Приложение к Блоку 1 — машинные якоря}
% =============================================================================
\small\raggedright\sloppy
\noindent Всё перечисленное машинно проверено в репозитории; единственная аксиома
ветви --- \texttt{classic} (= Закон Исключённого Третьего, L3). Строки сверяются при
финальной сборке.

\begin{itemize}\itemsep2pt
\item \textbf{§1.1.} Самосвидетельствование $\neg\exists\to\exists$ --- онтологический
довод (полное обоснование --- Том~I); формальная цепь по построению начинается со
следствия (акта различения), а не с самого $A=\exists$.

\item \textbf{§1.2, структура.} Запись \texttt{Distinction} с полями
\texttt{positive}, \texttt{negative}, \texttt{exclusive}~(L2), \texttt{exhaustive}~(L3);
со\-конституция (обе стороны из одного акта) --- \texttt{co\_constitution}; пять свойств
различения --- \texttt{five\_properties\_of\_distinction}.\\
\emph{Файл:} \texttt{src/foundation/Distinction.v}.

\item \textbf{§1.2, единица (ядро позитивной онтологии).} «Один акт = единица» ---
\texttt{one\_distinction\_exists} (счёт \texttt{1}, свидетель
\texttt{[distinction\_of True]}); «пустой набор = нуль = отсутствие» ---
\texttt{zero\_distinctions} (счёт \texttt{0}, свидетель \texttt{[]});
счёт как добавление --- \texttt{distinction\_count\_succ}.\\
\emph{Файл:} \texttt{src/foundation/Distinction.v}, секция \emph{UNIT}.

\item \textbf{§1.2, асимметрия.} $\mathit{positive}\,D\to\neg\,\mathit{negative}\,D$ ---
\texttt{L4\_self\_grounding}; в форме закона --- \texttt{Law\_of\_SufficientReason}.\\
\emph{Файлы:} \texttt{Distinction.v}, \texttt{LawsFromDistinction.v}.

\item \textbf{Триада E/R/R} (три грани одного акта) --- экспозиция в
\texttt{docs/ERR\_FRAMEWORK.md}; формализация триады --- §1.5 (запись
\texttt{FunctionalSystem}, файл \texttt{TheoryOfSystems\_Core\_ERR.v}, §XVII).
\item \textbf{§1.3.} Пять законов из различения --- капстоун
\texttt{five\_laws\_from\_distinction}; L1 \texttt{L1\_stability}, L2 \texttt{L2\_exclusivity},
L3 \texttt{L3\_totality}~(=\,\texttt{classic}), L4 \texttt{L4\_self\_grounding},
L5 \texttt{L5\_hierarchy}.\\
\emph{Файлы:} \texttt{Distinction.v}, \texttt{LawsFromDistinction.v}.

\item \textbf{§1.4.} P1 \texttt{P1\_from\_L1\_L5}/\texttt{P1\_blocks\_russell};
P2 \texttt{P2\_from\_L5}; P3 \texttt{P3\_from\_L1\_L4}; P4 (по типовой конструкции)
\texttt{P4\_from\_L5}/\texttt{P4\_no\_completed\_infinity}, конструктив свидетеля ---
\texttt{L4\_witness} (\texttt{ToS\_Axioms.v}); капстоун
\texttt{four\_principles\_from\_five\_laws}.\\
\emph{Файл:} \texttt{PrinciplesFromLaws.v}.

\item \textbf{§1.5.} Триада \texttt{FunctionalSystem}, экстракторы
\texttt{get\_Elements/Roles/Rules}; вывод E/R/R из различения ---
\texttt{ERRFromDistinction.v}; экспозиция --- \texttt{docs/ERR\_FRAMEWORK.md}.\\
\emph{Файл:} \texttt{TheoryOfSystems\_Core\_ERR.v}, §XVII.
\end{itemize}
\normalsize

% --- Блок 1 завершён (1.1–1.5). Дальше: Блоки 2–8 по тезисным планам в Книги/Теория Систем/. ---

\end{document}
